%%% BEGIN fig:appO-cutsum
\begin{figure}[htbp]
\centering
\begin{tikzpicture}[fd,
  ff chain/.style={draw=blue!55!black,fill=blue!12,line width=0.6pt},
  ff soft/.style={gluon,draw=green!45!black},
  ff ins/.style={circle,draw=green!40!black,fill=green!45!black,inner sep=0pt,minimum size=4.2pt},
  ff lab/.style={font=\footnotesize},
  ff note/.style={font=\footnotesize,text=black!70},
  ff w/.style={font=\footnotesize,text=black!80},
  ff ctp/.style={draw=black!70,line width=0.9pt,postaction={decorate,decoration={markings,
      mark=at position #1 with {\arrow{Stealth[length=5pt,width=4.5pt]}}}}},
  ff ctp/.default=0.5,
  ff unit/.style={fill=gray!16,draw=gray!55,line width=0.4pt,rounded corners=2pt}]
  % ================================================================ (a) cut-position identity
  \begin{scope}
    \node[font=\small,anchor=west] at (-0.55,2.3) {(a)\ 割位置恒等式（含插入）};
    % closed-time-path order of the chain
    \draw[lcaxis,gray!60!black] (0.35,1.5) -- (8.45,1.5);
    \node[ff note,anchor=south] at (4.4,1.52) {闭时路序：振幅一侧 $\tau$ 递增，然后共轭一侧 $\tau$ 递减};
    % chain of final states between V_H and V_H^dagger
    \fill[ff chain] (0.3,-0.2) rectangle (8.5,0.2);
    \filldraw[hard blob] (0,0) circle[radius=0.3];
    \filldraw[hard blob] (8.8,0) circle[radius=0.3];
    \node[ff lab,anchor=north] at (0,-0.33) {$V_H$};
    \node[ff lab,anchor=north] at (8.8,-0.33) {$V_H^\dagger$};
    \foreach \x/\v in {2.2/1,4.4/2,6.6/3} {
      \node[vertex] at (\x,0) {};
      \node[ff lab,anchor=south] at (\x,0.22) {$v_\v$};
    }
    % insertions: v_1 on the amplitude side, v_3 on the conjugate side
    \draw[ff soft] (2.2,-1.0) -- (2.2,-0.02);
    \draw[ff soft] (6.6,-1.0) -- (6.6,-0.02);
    \node[ff ins] at (2.2,0) {};  \node[ff ins] at (6.6,0) {};
    \node[ff lab,anchor=north] at (2.2,-1.03) {振幅插入 $O_a$};
    \node[ff lab,anchor=north] at (6.6,-1.03) {共轭插入 $O_b'$};
    \node[ff lab,anchor=north] at (4.4,-0.27) {内部顶点};
    % weights of the vertices
    \node[ff w,anchor=north] at (2.2,-1.48) {$\alpha_1$，$\beta_1=0$};
    \node[ff w,anchor=north] at (4.4,-1.48) {$\alpha_2=\beta_2$};
    \node[ff w,anchor=north] at (6.6,-1.48) {$\alpha_3=0$，$\beta_3$};
    % cut positions through the states s_0 ... s_3
    \foreach \x/\j in {1.1/0,3.3/1,5.5/2,7.7/3} {
      \draw[cut] (\x,-0.55) -- (\x,0.85);
      \node[font=\footnotesize,text=red!65!black,anchor=south] at (\x,0.86) {$j=\j$};
      \node[ff lab,anchor=south west,inner sep=1pt] at (\x+0.04,0.2) {$s_\j$};
    }
    % table of weights
    \node[anchor=west,font=\footnotesize] at (9.55,0.2) {%
      \renewcommand{\arraystretch}{1.12}%
      \begin{tabular}{@{}c@{\quad}l@{}}
        \toprule
        $j$ & 割线穿过 $s_j$ 的权重 $W_j$\\
        \midrule
        0 & $\beta_1\beta_2\beta_3=0$\\
        1 & $\alpha_1\beta_2\beta_3$\\
        2 & $\alpha_1\alpha_2\beta_3$\\
        3 & $\alpha_1\alpha_2\alpha_3=0$\\
        \bottomrule
      \end{tabular}};
    % result
    \node[anchor=north,font=\small,align=center] at (6.6,-2.05) {%
      $\displaystyle\sum_{j=0}^{3}\mathcal T_j=-\ii\,\alpha_1\beta_2\beta_3\!\int\!\frac{\dd\mathcal E}{2\pi}\,H_1(\mathcal E)
       +\ii\,\alpha_1\alpha_2\beta_3\!\int\!\frac{\dd\mathcal E}{2\pi}\,H_3(\mathcal E)$，\\[3pt]
      $\displaystyle H_m(\mathcal E)=\prod_{i<m}\frac{1}{\mathcal E-E_i+\ii0}\prod_{i\ge m}\frac{1}{\mathcal E-E_i-\ii0}$};
  \end{scope}
  % ================================================================ (b) closed-time-path form
  \begin{scope}[yshift=-6.3cm]
    \node[font=\small,anchor=west] at (-0.55,1.5) {(b)\ 闭时路形式};
    % part later than all insertions: U^dagger U = 1
    \fill[ff unit] (4.1,-1.8) rectangle (8.4,0.78);
    \node[ff note,anchor=south] at (6.25,0.12) {$\tau>\tau_{\max}$：$U^\dagger U=1$};
    % branches: amplitude (upper, T order) and conjugate (lower, anti-T order)
    \draw[ff ctp] (0,0) -- (7.7,0);
    \draw[ff ctp=0.72] (7.7,-1.2) -- (0,-1.2);
    \draw[draw=black!70,line width=0.9pt] (7.7,0) arc[start angle=90,end angle=-90,radius=0.6];
    \filldraw[hard blob] (0,0) circle[radius=0.22];
    \filldraw[hard blob] (0,-1.2) circle[radius=0.22];
    \node[ff lab,anchor=east] at (-0.28,0) {$V_H(0)$};
    \node[ff lab,anchor=east] at (-0.28,-1.2) {$V_H^\dagger(0)$};
    \node[ff note,anchor=north west] at (0.3,-0.05) {振幅（${\rm T}$ 序）};
    \node[ff note,anchor=south west] at (0.3,-1.15) {共轭振幅（$\bar{\rm T}$ 序）};
    % insertions
    \draw[ff soft] (2.2,0.95) -- (2.2,0.03);   \node[ff ins] at (2.2,0) {};
    \draw[ff soft] (4.0,-2.15) -- (4.0,-1.23); \node[ff ins] at (4.0,-1.2) {};
    \node[ff lab,anchor=west] at (2.32,0.75) {$O_a(\tau_a)$};
    \node[ff lab,anchor=west] at (4.12,-2.0) {$O_b'^\dagger(\tau_b')$};
    % internal final-state interactions of J_A later than all insertions (they cancel)
    \foreach \x in {5.2,6.6} { \node[vertex] at (\x,0) {}; \node[vertex] at (\x,-1.2) {}; }
    \node[ff note,anchor=center] at (5.9,-0.6) {$J_A$ 内部的末态相互作用};
    % cut = final time
    \draw[cut] (8.3,-1.75) -- (8.3,0.9);
    \node[font=\footnotesize,text=red!65!black,anchor=west,align=left] at (8.45,-0.6) {割线\\（最终时刻 $\tau_f$）};
    % time axis
    \draw[lcaxis,gray!60!black] (0,-2.6) -- (9.6,-2.6) node[below left=1pt and -2pt,font=\footnotesize] {$\tau$};
    \foreach \x/\t in {0/0,2.2/{\tau_a},4.0/{\tau_b'=\tau_{\max}},8.3/{\tau_f}} {
      \draw[gray!60!black,line width=0.5pt] (\x,-2.6) -- ++(0,0.08);
      \node[font=\footnotesize,below] at (\x,-2.6) {$\t$};
    }
    % operator form after the sum over cuts
    \node[anchor=north,font=\small] at (4.8,-3.25) {%
      对割线求和：\quad$\sum_C\;\longrightarrow\;\langle P_A|\,\bar{\rm T}\big\{V_H^\dagger(0)\textstyle\prod_bO_b'^\dagger(\tau_b')\big\}\,
      {\rm T}\big\{V_H(0)\textstyle\prod_aO_a(\tau_a)\big\}|P_A\rangle$};
  \end{scope}
\end{tikzpicture}
\caption{割图求和的幺正性（\cref{lem:appO-cut-position}、\cref{lem:appO-unitarity}）。
(a)~$J_A$ 末态部分的一条链按闭时路序排列，割线穿过态 $s_j$ 的项 $\mathcal T_j$ 带权重 $W_j=\prod_{i\le j}\alpha_i\prod_{i>j}\beta_i$：振幅插入只在振幅一侧（$\beta=0$），共轭插入只在共轭一侧（$\alpha=0$），内部顶点两侧权重相同（$\alpha=\beta$，\cref{lem:appJ-cut-hermiticity}）。
把割线移过最晚的顶点只改变该顶点的权重，故“$v_2$ 最晚”的贡献 $\propto\alpha_2-\beta_2$ 消失，对割位置求和后只剩插入顶点给出的至多两项（$H_m$ 即“$v_m$ 最晚”）。
(b)~闭时路形式：对割线求和后 $J_A$ 成为图下所示的矩阵元，晚于所有插入的 $J_A$ 内部相互作用在两支上合成 $U^\dagger U=1$，即同一喷注内部的末态相互作用精确相消；软线把 $\tau>0$ 的插入固定在各自一侧，它们无法在 $J_A$ 内配对，这正是 Glauber 钳夹的来源。}
\label{fig:appO-cutsum}
\end{figure}
%%% END fig:appO-cutsum

%%% BEGIN fig:appO-deformation
\begin{figure}[htbp]
\centering
\begin{tikzpicture}[fd,
  ff pole/.pic={\draw[line width=1pt,line cap=round] (-0.085,-0.085)--(0.085,0.085) (-0.085,0.085)--(0.085,-0.085);},
  ff spole/.pic={\draw[line width=1pt,line cap=round,green!40!black] (-0.085,-0.085)--(0.085,0.085) (-0.085,0.085)--(0.085,-0.085);},
  ff contour/.style={draw=blue!65!black,line width=1.15pt,postaction={decorate,decoration={markings,
      mark=at position #1 with {\arrow{Stealth[length=6pt,width=5pt]}}}}},
  ff contour/.default=0.5,
  ff ax/.style={lcaxis,gray!60!black},
  ff win/.style={draw=red!55!black,line width=0.6pt},
  ff lab/.style={font=\footnotesize},
  ff note/.style={font=\footnotesize,text=black!75},
  ff dim/.style={draw=gray!60!black,line width=0.45pt,{Stealth[length=3pt]}-{Stealth[length=3pt]}}]
  % ================================================================ (a) complex k^- plane (J_A)
  \begin{scope}
    \fill[red!7] (-1.4,-1.0) rectangle (1.4,1.3);
    \draw[ff ax] (-3.3,0) -- (3.55,0) node[above left=0pt and -2pt] {$\Real k^-$};
    \draw[ff ax] (0,-1.05) -- (0,1.6) node[left] {$\Imag k^-$};
    \draw[ff contour=0.5] (-3.2,0) -- (-2.35,0) .. controls (-1.9,0) and (-1.85,0.85) .. (-1.4,0.85) -- (1.4,0.85)
         .. controls (1.85,0.85) and (1.9,0) .. (2.35,0) -- (3.2,0);
    \pic at (-0.5,-0.17) {ff pole}; \pic at (0.12,-0.24) {ff pole}; \pic at (0.48,-0.13) {ff pole};
    \node[ff note,anchor=north,align=center] at (0.1,-0.36) {$J_A$ 初态极点};
    \draw[ff dim] (1.05,0.02) -- (1.05,0.83);
    \node[ff lab,anchor=east] at (1.0,0.45) {$\kappa|\mathbf k_T|$};
    % window edges and the soft line's own pole (here k^+ < 0: upper half-plane, outside the window)
    \foreach \x in {-2.35,2.35} \draw[ff win] (\x,-0.1) -- (\x,0.1);
    \node[ff note,anchor=north,text=red!55!black] at (2.35,-0.12) {窗口边缘};
    \pic at (-2.95,0.2) {ff spole};
    \node[ff note,anchor=south,text=green!35!black,align=center] at (-2.95,0.32) {软线极点\\（$k^+<0$）};
    \node[ff note,anchor=south west,text=red!55!black] at (0.12,1.3) {窗口核（$\psi=1$）};
    \node at (0,-1.95) {(a)\ $J_A$：复 $k^-$ 平面，围道上移};
  \end{scope}
  % ================================================================ (b) complex k^+ plane (J_B)
  \begin{scope}[xshift=7.45cm]
    \fill[red!7] (-1.4,-1.3) rectangle (1.4,1.0);
    \draw[ff ax] (-3.3,0) -- (3.55,0) node[above left=0pt and -2pt] {$\Real k^+$};
    \draw[ff ax] (0,-1.4) -- (0,1.3) node[left] {$\Imag k^+$};
    \draw[ff contour=0.5] (-3.2,0) -- (-2.35,0) .. controls (-1.9,0) and (-1.85,-0.85) .. (-1.4,-0.85) -- (1.4,-0.85)
         .. controls (1.85,-0.85) and (1.9,0) .. (2.35,0) -- (3.2,0);
    \pic at (-0.45,0.15) {ff pole}; \pic at (0.15,0.25) {ff pole}; \pic at (0.5,0.12) {ff pole};
    \node[ff note,anchor=south,align=center] at (0.05,0.35) {$J_B$ 初态极点};
    \draw[ff dim] (1.05,-0.02) -- (1.05,-0.83);
    \node[ff lab,anchor=east] at (1.0,-0.45) {$\kappa|\mathbf k_T|$};
    \foreach \x in {-2.35,2.35} \draw[ff win] (\x,-0.1) -- (\x,0.1);
    \node[ff note,anchor=south,text=red!55!black] at (-2.35,0.12) {窗口边缘};
    \pic at (2.95,-0.2) {ff spole};
    \node[ff note,anchor=north,text=green!35!black,align=center] at (2.95,-0.32) {软线极点\\（$k^->0$）};
    \node[ff note,anchor=north west,text=red!55!black] at (0.12,-1.3) {窗口核（$\psi=1$）};
    \node at (0,-1.95) {(b)\ $J_B$：复 $k^+$ 平面，围道下移};
  \end{scope}
  % ================================================================ (c) moduli plane at fixed |k_T| ~ lambda Q
  \begin{scope}[yshift=-4.45cm,xshift=1.6cm,x=1.32cm,y=1.32cm]
    % coordinates: ln|k^\pm|/ln(1/lambda);  -2 <-> lambda^2 Q, -1 <-> lambda Q, 0 <-> Q;  |k_T| = lambda Q.
    % lines |k^+||k^-| = const have slope -1; offsets from x+y=-2 schematic, drawn for kappa=1/2, delta=1/8:
    % c = min(kappa^2,delta/2) = delta/2 (the core boundary), kappa^2 = 1/4 between delta and 1/2.
    \fill[red!8] (-3.15,-3.15) -- (0.65,-3.15) -- (-3.15,0.65) -- cycle;                              % G_delta
    \draw[red!55!black,line width=0.75pt] (0.65,-3.15) -- (-3.15,0.65);                               % delta k_T^2
    \draw[red!55!black,line width=0.55pt,dash pattern=on 2.2pt off 1.5pt] (0.4,-3.15) -- (-3.15,0.4);  % delta k_T^2/2
    \draw[green!40!black,line width=0.75pt] (1.0,-2.85) -- (-2.85,1.0);                               % 2k^+k^- = k_T^2
    % axes and ticks
    \draw[ff ax] (-3.2,-3.2) -- (1.05,-3.2) node[below left=1pt and -4pt] {$\ln|k^+|$};
    \draw[ff ax] (-3.2,-3.2) -- (-3.2,1.05) node[left] {$\ln|k^-|$};
    \foreach \v/\t in {-2/{\lambda^2Q},-1/{\lambda Q},0/{Q}} {
      \draw[line width=0.5pt] (\v,-3.2) -- ++(0,0.06);  \node[font=\scriptsize,below] at (\v,-3.2) {$\t$};
      \draw[line width=0.5pt] (-3.2,\v) -- ++(0.06,0);  \node[font=\scriptsize,left] at (-3.2,\v) {$\t$};
    }
    % image of the window core under the deformation: |k^+|,|k^-| >= kappa |k_T|
    \fill[blue!10] (-1.09,-1.09) rectangle (0.75,0.75);
    \draw[blue!60!black,line width=0.6pt] (0.75,-1.09) -- (-1.09,-1.09) -- (-1.09,0.75);
    \node[circle,draw=red!60!black,line width=0.8pt,fill=white,inner sep=0pt,minimum size=6pt] (G) at (-2,-2) {};
    \node[circle,fill=blue!60!black,inner sep=0pt,minimum size=5pt] (D) at (-1.09,-1.09) {};
    \draw[blue!60!black,line width=0.75pt,-{Stealth[length=5pt,width=4pt]}] (G) -- (D);
    \node[font=\scriptsize,anchor=north,text=red!55!black,inner sep=1.5pt] at (G.south) {Glauber 实点};
    \node[font=\scriptsize,anchor=south west,text=blue!50!black,inner sep=1.5pt] at (-1.05,-1.05) {$|k^\pm|\ge\kappa|\mathbf k_T|$};
    \node[font=\scriptsize,text=green!35!black,anchor=south west,inner sep=1pt] at (0.05,-2.0) {软线质壳};
    \node[font=\small,text=red!55!black] at (-2.65,-2.75) {$G_\delta$};
    \node at (-1.05,-4.0) {(c)\ $|k^\pm|$ 的对数平面，$|\mathbf k_T|\sim\lambda Q$};
  \end{scope}
  % ---------------------------------------------------------------- notes: constants, line styles, eikonal bound
  \node[anchor=north west,font=\footnotesize,align=left,text width=7.1cm] at (3.55,-3.3) {%
    形变：$k\to k+\ii\sigma\kappa\psi|\mathbf k_T|\,e_G$，$e_G=n_B-n_A$。\\[2pt]
    (c) 中斜线（$|k^+||k^-|=$ 常数）：红实线 $\delta\mathbf k_T^2$（$G_\delta$ 的边界，形变在其外关闭），红虚线 $\frac12\delta\mathbf k_T^2$（窗口核边界），绿线为连接线的质壳 $2k^+k^-=\mathbf k_T^2$；图中取 $\kappa=\frac12$、$\delta=\frac18$，此时 $c=\delta/2$，红虚线即 $c\,\mathbf k_T^2$。\\[2pt]
    \textbullet\ 窗口核上 $|k^+||k^-|\ge|{\rm Im}\,k^+||{\rm Im}\,k^-|=\kappa^2\mathbf k_T^2$，核外实部已满足 $|k^+k^-|>\frac12\delta\mathbf k_T^2$，故形变后处处 $|k^+||k^-|\ge c\,\mathbf k_T^2$，$c=\min(\kappa^2,\delta/2)$；\\[1pt]
    \textbullet\ $\mathcal G_0$：$0<\kappa\le\frac12$，$0<\delta\le\frac18$；分级构型另需 $\kappa^2<1/(2KN_C^2)$；\\[1pt]
    \textbullet\ 形变后的连接线 ${\rm Re}(k^2)\le-\frac14\mathbf k_T^2$，不触及质壳；\\[1pt]
    \textbullet\ 实围道上程函近似以 $O(1)$ 失效；形变后的围道上 Grammer--Yennie 比值为 $O(\lambda/\kappa)$。};
\end{tikzpicture}
\caption{单侧改写之后的同时形变（振幅一侧 $\sigma=+1$，软动量 $k$ 流入 $J_A$；共轭一侧方向相反；示意，非比例）。
(a)~改写后 $J_A$ 只经初态分母依赖 $k^-$，其极点全在下半平面；围道在窗口核上抬到 ${\rm Im}\,k^-=\kappa|\mathbf k_T|$，在窗口边缘之前回到实轴，因而不触及窗口外的连接线自身极点 $k^-=(\mathbf k_T^2-\ii0)/(2k^+)$（$k^+>0$ 时它在下半平面，不受影响）。
(b)~$J_B$ 的 $k^+$ 极点全在上半平面，围道同时下移到 ${\rm Im}\,k^+=-\kappa|\mathbf k_T|$，故 ${\rm Im}\,k\propto e_G=n_B-n_A$；只形变一个分量时 $|k^+k^-|$ 仍远小于 $\mathbf k_T^2$。
(c)~以 $\ln|k^\pm|$ 为坐标（$|\mathbf k_T|\sim\lambda Q$）：窗口核中的 Glauber 点（$|k^\pm|\sim\lambda^2Q$，位于 $G_\delta=\{|k^+k^-|\le\delta\mathbf k_T^2\}$ 深处）被移到 $|k^\pm|\ge\kappa|\mathbf k_T|$ 处，形变后的围道处处满足 $|k^+||k^-|\ge c\,\mathbf k_T^2$，从而离开 Glauber 区，程函近似在其上有定量误差界（\cref{lem:appO-no-crossing}、\cref{prop:appO-target}、\cref{prop:appO-eikonal}）。}
\label{fig:appO-deformation}
\end{figure}
%%% END fig:appO-deformation

%%% BEGIN fig:appO-flowchart
\begin{figure}[htbp]
\centering
\begin{tikzpicture}[fd,
  ff box/.style={draw=black!55,line width=0.6pt,rounded corners=2pt,fill=gray!3},
  ff title/.style={font=\footnotesize\bfseries,inner sep=0pt,anchor=north west},
  ff err/.style={font=\scriptsize,text=black!65,inner sep=0pt,anchor=north east},
  ff key/.style={font=\scriptsize,text=black!75,inner sep=0pt,anchor=north west},
  ff row/.style={font=\scriptsize,text=black!60,inner sep=0pt,anchor=west},
  ff tag/.style={draw,line width=0.45pt,rounded corners=1pt,inner xsep=1.8pt,inner ysep=0pt,minimum height=3.7mm,
                 font=\scriptsize,anchor=west},
  ff R/.style={ff tag,fill=OliveGreen!15,draw=OliveGreen!65!black},
  ff Rc/.style={ff R,dash pattern=on 1.6pt off 1pt},
  ff A/.style={ff tag,fill=RoyalBlue!15,draw=RoyalBlue!65!black},
  ff M/.style={ff tag,fill=white,draw=black!60},
  ff D/.style={ff tag,fill=none,draw=black!45,line width=0.35pt},
  ff asm/.style={draw=orange!65!black,fill=orange!8,line width=0.6pt,dash pattern=on 2.4pt off 1.4pt,rounded corners=1pt,
                 font=\scriptsize,align=left,text width=5.0cm,inner sep=2.2pt,anchor=north west},
  ff asmN/.style={ff asm,draw=gray!60!black,fill=gray!7},
  ff badge/.style={draw=orange!65!black,fill=orange!8,line width=0.6pt,dash pattern=on 2.4pt off 1.4pt,rounded corners=1pt,
                   font=\scriptsize,inner sep=2.2pt,anchor=west},
  ff appN/.style={draw=OliveGreen!65!black,fill=OliveGreen!8,line width=0.6pt,rounded corners=2pt,font=\scriptsize,
                  align=center,inner sep=3pt},
  ff flow/.style={draw=black!65,line width=0.75pt,-{Stealth[length=4.2pt,width=3.6pt]}},
  ff use/.style={draw=orange!65!black,line width=0.55pt,dash pattern=on 2.4pt off 1.4pt,-{Stealth[length=3.6pt,width=3pt]}},
  ff useN/.style={ff use,draw=gray!55!black}]
  % ------------------------------------------------------------------ geometry
  % step boxes: x in [0,9.9]; assumption column: x in [10.55,15.8]
  \def\ffR{9.9}
  % ================================================================== (O-1)
  \draw[ff box] (0,0) rectangle (\ffR,-1.33);
  \node[ff title] at (0.15,-0.12) {(O-1)\quad 区域、积分域与窗口};
  \node[ff err] at (\ffR-0.15,-0.13) {误差：无};
  \node[ff D] (a1) at (0.15,-0.66) {\cref{def:appO-leading-data}};
  \node[ff D,right=1.2mm of a1] (a2) {\cref{def:appO-windows}};
  \node[ff D,right=1.2mm of a2] (a3) {\cref{def:appO-weak-equivalence}};
  \node[ff key] at (0.15,-0.98) {领头区域 $R$ 与支配集 $U_R$；类 $\mathcal G_0$；窗口 $\Theta_j$ 与软角区 $K_R$；弱等价 $\doteq_R$};
  % ================================================================== (O-2)
  \begin{scope}[yshift=-1.83cm]
  \draw[ff box] (0,0) rectangle (\ffR,-1.33);
  \node[ff title] at (0.15,-0.12) {(O-2)\quad 光锥序微扰论中的领头幂次形式};
  \node[ff err] at (\ffR-0.15,-0.13) {误差：E-S、E-P；其余精确};
  \node[ff R] (b1) at (0.15,-0.66) {\cref{lem:appO-lcopt}};
  \node[ff R,right=1.2mm of b1] (b2) {\cref{lem:appO-numerators}};
  \node[ff R,right=1.2mm of b2] (b3) {\cref{lem:appO-orientation}};
  \node[ff R,right=1.2mm of b3] (b4) {\cref{lem:appO-glauber-pc}};
  \node[ff Rc,right=1.2mm of b4] (b5) {\cref{lem:appO-hard-locality}};
  \node[ff key] at (0.15,-0.98) {$J_A$：时间 $\tau=x^+$，能量 $l^-$；$J_B$：时间 $x^-$，能量 $l^+$；$V_H$ 位于 $\tau=0$};
  \end{scope}
  % ================================================================== (O-3)
  \begin{scope}[yshift=-3.66cm]
  \draw[ff box] (0,0) rectangle (\ffR,-2.86);
  \node[ff title] at (0.15,-0.12) {(O-3)\quad 割图求和与单侧改写};
  \node[ff err] at (\ffR-0.15,-0.13) {误差：精确；E-F；E-W};
  \node[ff row] at (0.15,-0.66) {割图求和};
  \node[ff M] (c1) at (1.75,-0.66) {\cref{lem:appO-cut-position}};
  \node[ff Rc,right=1.2mm of c1] (c2) {\cref{lem:appO-unitarity}};
  \node[ff row] at (0.15,-1.13) {等时局域化};
  \node[ff M] (c3) at (1.75,-1.13) {\cref{lem:appO-equal-time}};
  \node[ff R,right=1.2mm of c3] (c4) {\cref{lem:appO-insertion-values}};
  \node[ff Rc,right=1.2mm of c4] (c5) {\cref{prop:appO-ss-localize}};
  \node[ff R,right=1.2mm of c5] (c6) {\cref{lem:appO-joint-localization}};
  \node[ff row] at (0.15,-1.60) {旁观者交换};
  \node[ff R] (c7) at (1.75,-1.60) {\cref{lem:appO-glauber-operator}};
  \node[ff Rc,right=1.2mm of c7] (c8) {\cref{prop:appO-ss-cancel}};
  \node[ff row] at (0.15,-2.07) {单侧改写};
  \node[ff M] (c9) at (1.75,-2.07) {\cref{lem:appO-carrier}};
  \node[ff D,right=1.2mm of c9] (c10) {\cref{def:appO-rewriting}};
  \node[ff Rc,right=1.2mm of c10] (c11) {\cref{prop:appO-one-sided}};
  \node[ff key] at (0.15,-2.42) {$\sum_C$：最大时刻配对；插入局域于 $\tau=0$；整函数载子 $\phi(x)=2\sqrt\pi\,\ee^{-x^2}$};
  \end{scope}
  % ================================================================== (O-4)
  \begin{scope}[yshift=-7.02cm]
  \draw[ff box] (0,0) rectangle (\ffR,-1.33);
  \node[ff title] at (0.15,-0.12) {(O-4)\quad 减除（链）项与改写相容};
  \node[ff err] at (\ffR-0.15,-0.13) {误差：E-W};
  \node[ff D] (d1) at (0.15,-0.66) {\cref{def:appO-rewriting}};
  \node[ff Rc,right=1.2mm of d1] (d2) {\cref{prop:appO-one-sided}};
  \node[font=\scriptsize,text=black!60,anchor=west,right=2mm of d2,inner sep=0pt] {（链项部分；需 \cref{asm:appN-separation}(ii)）};
  \node[ff key] at (0.15,-0.98) {先 $T_c$、后改写：$C_{R'}\tilde F\equiv\mathfrak r_RC_{R'}F_\Gamma$；(ET1)(ET2) 保持，$\lambda_J$ 不增};
  \end{scope}
  % ================================================================== (O-5)
  \begin{scope}[yshift=-8.85cm]
  \draw[ff box] (0,0) rectangle (\ffR,-1.33);
  \node[ff title] at (0.15,-0.12) {(O-5)\quad 同时形变：形变目标};
  \node[ff err] at (\ffR-0.15,-0.13) {误差：精确（几何）};
  \node[ff D] (e1) at (0.15,-0.66) {\cref{def:appO-deformation-map}};
  \node[ff R,right=1.2mm of e1] (e2) {\cref{lem:appO-no-crossing}};
  \node[ff D,right=1.2mm of e2] (e3) {\cref{def:appO-graded}};
  \node[ff D,right=1.2mm of e3] (e4) {\cref{def:appO-target}};
  \node[ff R,right=1.2mm of e4] (e5) {\cref{prop:appO-target}};
  \node[ff key] at (0.15,-0.98) {$k^-\to k^-+\ii\sigma\kappa\psi|\mathbf k_T|$，$k^+\to k^+-\ii\sigma\kappa\psi|\mathbf k_T|$（直接线，沿 $e_G=n_B-n_A$）};
  \end{scope}
  % ================================================================== (O-6)
  \begin{scope}[yshift=-10.68cm]
  \draw[ff box] (0,0) rectangle (\ffR,-1.33);
  \node[ff title] at (0.15,-0.12) {(O-6)\quad 程函近似与指向过去的 Wilson 线};
  \node[ff err] at (\ffR-0.15,-0.13) {误差：E-P（于 $\Omega$ 上）};
  \node[ff R] (f1) at (0.15,-0.66) {\cref{prop:appO-eikonal}};
  \node[ff A,right=1.2mm of f1] (f2) {\cref{prop:appO-past-wilson}};
  \node[font=\scriptsize,text=black!60,anchor=west,right=2mm of f2,inner sep=0pt] {（需 \cref{asm:appN-ward-product}）};
  \node[ff key] at (0.15,-0.98) {程函分母 $n_A\cdot K+\ii\sigma0$：$W[0,-\infty n_A]$、$W[0,-\infty n_B]$（指向过去）};
  \end{scope}
  % ================================================================== (O-7)
  \begin{scope}[yshift=-12.51cm]
  \draw[ff box] (0,0) rectangle (\ffR,-1.33);
  \node[ff title] at (0.15,-0.12) {(O-7)\quad 软因子 $=1$，共线因子 $=$ DIS 的 PDF};
  \node[ff err] at (\ffR-0.15,-0.13) {误差：精确};
  \node[ff R] (g1) at (0.15,-0.66) {\cref{lem:method-soft-unitarity}};
  \node[ff key] at (0.15,-0.98) {$S^{(-)}$ 由指向过去的 Wilson 线构成，对软末态求和后 $\mathcal U^{(-)\dagger}\mathcal U^{(-)}=1$};
  \end{scope}
  % ================================================================== (O-8)
  \begin{scope}[yshift=-14.34cm]
  \draw[ff box] (0,0) rectangle (\ffR,-1.33);
  \node[ff title] at (0.15,-0.12) {(O-8)\quad 主定理与附录 N 的接口};
  \node[ff err] at (\ffR-0.15,-0.13) {误差：E-W};
  \node[ff Rc] (h1) at (0.15,-0.66) {\cref{thm:appO-main}};
  \node[ff Rc,right=1.2mm of h1] (h2) {\cref{cor:appO-interface}};
  \node[ff Rc,right=1.2mm of h2] (h3) {\cref{prop:appO-tmd}};
  \node[ff key] at (0.15,-0.98) {(Gla-a)--(Gla-c)：弱等价、单侧取向（$\eta=-1$）、形变目标};
  \end{scope}
  % ------------------------------------------------------------------ step arrows
  \foreach \ya/\yb in {-1.33/-1.83, -3.16/-3.66, -6.52/-7.02, -8.35/-8.85, -10.18/-10.68, -12.01/-12.51, -13.84/-14.34}
    \draw[ff flow] (4.95,\ya) -- (4.95,\yb);
  % ------------------------------------------------------------------ App N (green) and the users
  \node[ff appN,anchor=north] (N) at (4.95,-16.17)
    {附录 N：\cref{prop:appN-dy}（DY 的减除定理）、\cref{prop:appN-tmd-window}（TMD 窗口）};
  \draw[ff flow] (4.95,-15.67) -- (N.north);
  \node[font=\scriptsize,anchor=north west,inner sep=0pt] (U0) at (0.15,-16.98) {使用者：};
  \node[ff A,right=1mm of U0] (U1) {\cref{thm:dy-factorization}};
  \node[ff A,right=1.2mm of U1] (U2) {\cref{lem:tmd-glauber}};
  \node[ff A,right=1.2mm of U2] (U3) {\cref{prop:method-glauber-deformation}};
  \node[font=\scriptsize,text=black!65,anchor=west,right=1.5mm of U3,inner sep=0pt] {（第 12、14、9 章，全阶论证 [A]）};
  % ------------------------------------------------------------------ assumptions (orange, dashed)
  \node[ff asm] (L) at (10.55,-2.25) {\cref{asm:appO-lcopt}\ 光锥序微扰论的适用性：\\无零模与端点贡献；瞬时圈取值，(i)(ii)};
  \node[ff asm] (X) at (10.55,-4.30) {\cref{asm:appO-mixing}\ Glauber--软混合与重合连接：\\(a)(b) 一般软子图；(c) 完整图集相消};
  \node[ff asm] (V) at (10.55,-5.70) {\cref{asm:appO-envelope}\ 纤维积分误差的可积包络};
  \node[ff asmN] (NN) at (10.55,-6.75) {附录 N 的共同输入（与 DIS、$e^+e^-$\\共用）：\cref{asm:appN-envelope}、\cref{asm:appN-separation}、\\\cref{asm:appN-ward-product}};
  \node[ff asm] (F) at (10.55,-8.75) {\cref{asm:appO-deformation}\ (a) 一般构型的形变；\\(b) 与集代数 $\mathfrak B$ 相容；(c) 软角区；\\(d) H-发射端点与重合 SA/AS 插入};
  \node[ff asm] (C) at (10.55,-10.85) {\cref{asm:appO-rewritten-cs}\ 改写形式与族求和形式的\\(CS1)(CS2)（含处方的逐项匹配）};
  % uses: assumption -> step (east side of the step boxes)
  \draw[ff use] (L.west) -- (\ffR,-2.50);
  \draw[ff use] (L.west) -- (\ffR,-3.95);
  \draw[ff use] (X.west) -- (\ffR,-4.70);
  \draw[ff use] (V.west) -- (\ffR,-5.85);
  \draw[ff use] (V.west) -- (\ffR,-7.40);
  \draw[ff useN] (NN.west) -- (\ffR,-6.20);
  \draw[ff useN] (NN.west) -- (\ffR,-7.70);
  \draw[ff use] (F.west) -- (\ffR,-9.50);
  \draw[ff use] (C.west) -- (\ffR,-11.50);
  % (O-8) uses asm:appO-deformation (b)(c) and asm:appO-rewritten-cs again: short badges
  \node[ff badge] (F8) at (10.55,-14.72) {\cref{asm:appO-deformation}(b)(c)};
  \node[ff badge] (C8) at (10.55,-15.22) {\cref{asm:appO-rewritten-cs}};
  \draw[ff use] (F8.west) -- (\ffR,-14.72);
  \draw[ff use] (C8.west) -- (\ffR,-15.22);
  % ------------------------------------------------------------------ legend
  \begin{scope}[yshift=-17.75cm]
    \node[ff R] (l1) at (0,0) {\phantom{O}};   \node[font=\scriptsize,anchor=west,right=0.8mm of l1,inner sep=0pt] (m1) {无条件证明（$R_{\rm PT}$/$R_{\rm def}$）};
    \node[ff Rc,right=3mm of m1] (l2) {\phantom{O}}; \node[font=\scriptsize,anchor=west,right=0.8mm of l2,inner sep=0pt] (m2) {在编号假设下证明（$R_{\rm cond}$）};
    \node[ff M,right=3mm of m2] (l3) {\phantom{O}}; \node[font=\scriptsize,anchor=west,right=0.8mm of l3,inner sep=0pt] (m3) {数学引理};
    \node[ff A,right=3mm of m3] (l4) {\phantom{O}}; \node[font=\scriptsize,anchor=west,right=0.8mm of l4,inner sep=0pt] (m4) {全阶论证 [A]};
    \node[ff D,right=3mm of m4] (l5) {\phantom{O}}; \node[font=\scriptsize,anchor=west,right=0.8mm of l5,inner sep=0pt] (m5) {定义};
    \node[ff tag,draw=orange!65!black,fill=orange!8,line width=0.6pt,dash pattern=on 2.4pt off 1.4pt,right=3mm of m5] (l6) {\phantom{O}}; \node[font=\scriptsize,anchor=west,right=0.8mm of l6,inner sep=0pt] {编号假设};
  \end{scope}
\end{tikzpicture}
\caption{附录~O 中单举色单态 Drell--Yan 的 Glauber 区相消条件证明的结构：步骤 (O-1)--(O-8) 自上而下，每个框列出建立该步的陈述（框的样式同附录~P 的依赖图，见图例）、该步的误差类型（E-S 标度陈述、E-P 逐点界、E-F 纤维积分界、E-W 弱等价）与要点。
橙色虚线框为附录~O 的编号假设，虚线箭头指向使用它们的步骤；灰色虚线框为与 DIS、$e^+e^-$ 共用的附录~N 输入。
\cref{asm:appO-deformation} 的 (b)(c)(d) 即使对阿贝尔理论的类 $\mathcal G_0$ 也需要，因此在截面层次上 DY 的 Glauber 相消没有一步是无条件的。
步骤 (O-8) 经 \cref{cor:appO-interface} 与 \cref{prop:appO-tmd} 把结果交给附录~N 的减除定理，供第~12、14、9~章的全阶论证使用。}
\label{fig:appO-flowchart}
\end{figure}
%%% END fig:appO-flowchart

%%% BEGIN fig:appO-lcopt-term
\begin{figure}[htbp]
\centering
\begin{tikzpicture}[fd,
  ff hadron/.style={draw=black!70,line width=2.2pt,postaction={decorate,decoration={markings,
      mark=at position 0.55 with {\arrow[black!70]{Stealth[length=7pt,width=7.5pt]}}}}},
  ff soft/.style={gluon,draw=green!45!black},
  ff state/.style={draw=gray!60!black,line width=0.55pt,dash pattern=on 2.2pt off 1.6pt},
  ff ax/.style={lcaxis,gray!60!black},
  ff lab/.style={font=\footnotesize},
  ff dlab/.style={font=\small,fill=white,inner sep=1.2pt},
  ff note/.style={font=\footnotesize,text=black!70}]
  % ---------------------------------------------------------------- light-cone times
  \def\tA{-5.0}\def\tB{-2.4}\def\tC{2.2}\def\tD{4.3}\def\tcut{5.6}
  % initial / final regions (light shading)
  \fill[blue!4] (\tA,-1.2) rectangle (0,3.2);
  \fill[orange!6] (0,-1.2) rectangle (\tcut,3.2);
  \node[ff note,anchor=south] at ({(\tA)/2},3.2) {初态（$\tau<0$）};
  \node[ff note,anchor=south] at ({(\tcut)/2},3.2) {末态（$\tau>0$，割线之前）};
  % ---------------------------------------------------------------- lines of J_A (amplitude side)
  % hadron A -> contact vertex (v1, v1') at tau_1; instantaneous quark line v1 -> v1'
  \coordinate (v1) at (\tA,0);  \coordinate (v1p) at (\tA,1.3);
  \coordinate (VH) at (0,1.3);
  \coordinate (w2) at (\tB,0);  \coordinate (w3) at (\tC,0);
  \coordinate (v4) at (\tD,2.35);
  \draw[ff hadron] (-6.75,0) -- (v1);
  \draw[quark=0.86] (v1) -- (v1p);
  \draw[line width=0.9pt] ($(v1)!0.42!(v1p)+(-0.16,0)$) -- ++(0.32,0);   % bar: instantaneous part
  \draw[quark] (v1p) -- (VH);                                          % active quark a' -> V_H
  \draw[gluon] (v1p) to[out=22,in=185] (v4);                           % collinear gluon g
  \draw[scalar] (v1) -- (w2) -- (w3) -- (\tcut+0.6,0);                 % scalar spectator s, s', s''
  \draw[quark] (v4) -- (\tcut+0.6,2.85);                               % final-state splitting g -> q qbar
  \draw[quark] (\tcut+0.6,1.85) -- (v4);
  % soft injections k_1 (initial) and k_2 (final), oriented into J_A
  \draw[ff soft] (\tB,-1.1) -- (w2);
  \draw[ff soft] (\tC,-1.1) -- (w3);
  \path[mom={$k_1$},mom sep=0.32cm] (\tB,-1.1) -- (w2);
  \path[mom'={$k_2$},mom sep=0.32cm] (\tC,-1.1) -- (w3);
  % vertices
  \foreach \p in {v1,v1p,w2,w3,v4} \node[vertex] at (\p) {};
  \filldraw[hard blob] (VH) circle[radius=0.2];
  % labels of lines
  \node[ff lab,anchor=north] at (-6.1,-0.12) {$P_A$};
  \node[ff lab,anchor=east] at ($(v1)!0.42!(v1p)+(-0.22,0)$) {瞬时线};
  \node[ff lab,anchor=north] at ($(v1p)!0.5!(VH)+(0,-0.06)$) {$a'$};
  \node[ff lab,anchor=south] at (-4.0,2.05) {$g$};
  \node[ff lab,anchor=south] at (-4.4,0.04) {$s$};
  \node[ff lab,anchor=south] at (-0.6,0.04) {$s'$};
  \node[ff lab,anchor=south] at (4.85,0.04) {$s''$};
  \node[ff lab,anchor=west] at (\tcut+0.6,2.85) {$q$};
  \node[ff lab,anchor=west] at (\tcut+0.6,1.85) {$\bar q$};
  \node[ff lab,anchor=north west] at ($(VH)+(0.1,-0.14)$) {$V_H$};
  % ---------------------------------------------------------------- intermediate states and the cut
  \foreach \x/\d in {-3.7/1,-1.2/2,1.1/3,3.25/4} {
    \draw[ff state] (\x,-1.15) -- (\x,2.75);
    \node[ff dlab] at (\x,2.95) {$D_\d$};
  }
  \draw[cut] (\tcut,-1.15) -- (\tcut,3.2);
  \node[font=\footnotesize,text=red!65!black,anchor=north west] at (\tcut+0.05,-0.2) {割线};
  % ---------------------------------------------------------------- time axis
  \draw[ff ax] (-6.75,-1.5) -- (6.7,-1.5) node[below left=1pt and -2pt,font=\footnotesize] {$\tau=x^+$};
  \foreach \x/\t in {\tA/{\tau_1},\tB/{\tau_2},0/{0},\tC/{\tau_3},\tD/{\tau_4}} {
    \draw[gray!60!black,line width=0.5pt] (\x,-1.5) -- ++(0,0.08);
    \node[font=\footnotesize,below] at (\x,-1.5) {$\t$};
  }
  % ---------------------------------------------------------------- energy denominators
  \node[anchor=north,font=\small,align=left] at (0,-2.1) {%
    $\begin{array}{@{}l@{\qquad}l@{}}
      D_1=P_A^--E_s-E_{a'}-E_g+\ii0, & D_3=\mathcal E-k_2^--E_{s'}-E_g+\ii0,\\[2pt]
      D_2=P_A^-+k_1^--E_{s'}-E_{a'}-E_g+\ii0, & D_4=\mathcal E-E_{s''}-E_g+\ii0,\\[2pt]
      \multicolumn{2}{@{}l@{}}{\text{割线：}2\pi\delta\big(\mathcal E-E_{s''}-E_q-E_{\bar q}\big),\qquad
        E_l\equiv E(l)=\dfrac{\mathbf l_T^2+m_l^2-\ii0}{2l^+}.}
    \end{array}$};
\end{tikzpicture}
\caption{一个光锥序微扰论项（$J_A$ 的振幅一侧，$\mathcal G_0$ 的旁观者模型：强子 $A$ 分裂为活动夸克与标量旁观者 $s$；共轭一侧镜像，未画出）。
顶点按光锥时间 $\tau=x^+$ 排序，相邻顶点之间的每个中间态（灰色虚线）给出一个能量分母 $\ii/D_i$，$D_i$ 为注入的负分量 $l^-$ 减去该态各线的在壳能量 $E(l)$，再加 $+\ii0$（\cref{lem:appO-lcopt}）；
带短横的竖线是 Feynman 规范中费米子传播子的瞬时部分 $\ii\gamma^+/(2l^+)$，它把两端顶点置于同一时刻（\cref{lem:appO-numerators}）；硬顶点 $V_H$（与 $H$、$J_B$ 相连，未画出）位于 $\tau=0$。
软动量 $k_1$、$k_2$ 一律取为流入 $J_A$：$\tau<0$ 的初态分母形如 $A_D+k^-+\ii0$，零点在 ${\rm Im}\,k^-<0$；$\tau>0$ 的末态分母用末态总能量 $\mathcal E$ 写成 $A'_D-k^-+\ii0$，零点在 ${\rm Im}\,k^->0$（\cref{lem:appO-orientation}），
故 $k_1$ 只出现在 $D_2$，$k_2$ 只出现在 $D_3$。}
\label{fig:appO-lcopt-term}
\end{figure}
%%% END fig:appO-lcopt-term
